\chapter{Proposed Work, Deliverables, and Feasibility}
\label{ch:plan}
\section{Work Packages}
\subsection{Reference and Evaluation Control}
First, preserve the present data as the developmental baseline and create
new fixed evaluation points and starting-position sets. Check FEM mesh and
geometric refinement beyond the present pair where feasible. Repeat selected
particle paths with tighter ODE tolerances, smaller maximum steps, and
refined wall-event localization. This separates trajectory-integrator
sensitivity from flow-surrogate error before retraining is evaluated.

\subsection{Transport-Directed PINN Improvement}
Compare the current physics-and-flux loss with controlled alternatives:
increased boundary sampling, targeted collocation near turning regions,
and residual-informed sample redistribution. Use equal documented training
budgets where possible and record any extra compute. Develop each change
on the development set, then freeze the procedure before evaluating held-out
locations. The primary endpoints will be path error, time coverage and
endpoint agreement, in addition to velocity and pressure norms.

The main force ablation will retain exact Stokes drag on the learned flow.
It prevents a flow improvement from being confused with a change in the
force surrogate. If a force network provides no measurable benefit for
this known algebraic law, exact drag may remain the preferred computational
baseline. A future empirical force model would need external force data and
a separate identification/validation protocol.

\subsection{Geometry and Parameter Coverage}
Use straight, single-bend, S-shaped and smooth custom-centerline channels
at a modest set of Reynolds numbers, initially $0.5$, $1$ and $3$. These
are proposed cases, not completed new results. Retrain each case first to
establish a reliable per-case baseline. Only a later model trained across
multiple parameter cases could support an unseen-parameter test. Held-out
centerline shapes and Reynolds numbers must be specified before any
parametric-generalization claim is made.

Retain the existing particle-size and density sensitivities, add starting
points near observed event transitions, and distinguish low-confinement
quantitative cases from larger-particle exploratory cases. Review
$\Rey_p$, diameter-to-width ratio and loading assumptions in each case.
Changes to force physics require new references, not simply more network
training.

\subsection{Three-Dimensional and Physical Extensions}
Improve pressure recovery on the ABC verification problem before extending
to a finite-depth channel. A 3D channel study will require a mesh with
sidewalls, a three-component velocity network, physically specified 3D
inlet/outlet conditions and an independently refined 3D CFD reference.
The planar extrusion will continue to be identified as a visual aid.

Physical validation is conditional on an accessible, documented measurement
or benchmark dataset with compatible geometry, operating conditions,
particle properties and contact behavior. If such data remain unavailable,
the final claim will be numerical verification and model-to-model agreement,
with experimental validation explicitly outstanding.

\section{Schedule and Milestones}
The following relative schedule begins after project approval; it does not
imply an institutional deadline or an approved research start date.
\begin{table}[!htbp]
\centering\small
\caption{Proposed six-month research schedule.}
\label{tab:schedule}
\begin{tabularx}{\textwidth}{@{}lYY@{}}
\toprule Period & Main activity & Reviewable milestone\\\midrule
Month 1 & Freeze development/evaluation separation; refine references and ODE events & Reference and temporal sensitivity record\\
Month 2 & Boundary and turning-region training ablations & Controlled field and trajectory comparison\\
Month 3 & Expand starting positions, particle cases and force isolation & Transport-error and endpoint-transition maps\\
Month 4 & Separately trained geometry/Reynolds-number study & Case matrix with explicit supported scope\\
Month 5 & Improve ABC pressure; define a bounded 3D channel extension & 3D verification assessment and feasibility decision\\
Month 6 & Consolidate uncertainty, reproducibility and available physical evidence & Final analysis, thesis/proposal updates and research manuscript outline\\
\bottomrule
\end{tabularx}
\end{table}

\section{Resources, Risks, and Decision Rules}
The current workflow runs on CPU with bounded training and integration
budgets. GPU use may reduce training time, but exact speedups and total
future compute requirements have not been benchmarked. Separate costs will
be recorded for training, reference assembly/solve, trajectory evaluation
and reporting; reloading weights is not a new training-time measurement.

\begin{table}[!htbp]
\centering\small
\caption{Principal research risks and responses.}
\label{tab:risks}
\begin{tabularx}{\textwidth}{@{}YY@{}}
\toprule Risk & Response and decision rule\\\midrule
Low global error but wrong particle event & Add local and event-based diagnostics; withhold a transport-accuracy claim until path criteria are met.\\
Reference or event error dominates & Refine the reference and time/event resolution before attributing the discrepancy to the PINN.\\
Unstable or costly training & Compare fixed budgets and multiple seeds; retain failed cases and report supported operating ranges.\\
Particle-model assumptions fail & Restrict quantitative claims or introduce and independently validate additional physics.\\
Experimental evidence is unavailable & Limit the final validation statement to numerical comparisons.\\
3D pressure remains inaccurate & Treat the benchmark as unresolved and defer broad 3D channel claims.\\
\bottomrule
\end{tabularx}
\end{table}

\section{Expected Contributions and Deliverables}
The proposed deliverables are a validated numerical reference hierarchy,
documented relationships between field and transport error, a controlled
geometry/parameter evaluation, and a reproducible release containing
configuration, trained weights, raw data, figures and a data-driven LaTeX
report. These are research targets whose success will be judged using the
frozen evaluation design. The present proposal already demonstrates the
reporting mechanism and identifies where the preliminary model needs
improvement; it does not assume that every proposed target will succeed.

No human or animal data are used in the present numerical work. Any later
experimental collaboration, data reuse, institutional submission or required
authorship/assistance disclosure will follow the applicable institutional
and dataset terms. The current repository does not establish a separate
third-party data license or an experimental approval.
